\section{Introduction}\label{sec:introduction}

\subsection{First-Principles Derivations}\label{sec:intro-first-principles}

This paper studies a simple question with a hard combinatorial core: when a system exposes many state coordinates, which of those coordinates are actually needed to preserve the optimal decision? The aim is not to list all observable state variables, but to isolate the smallest part of the state that still determines what the best action is.

We model this with a decision problem: $A$ is the finite set of possible actions, $S$ is the state space, $U(a,s)$ is the utility of taking action $a$ in state $s$, and $\Opt(s)$ is the set of actions that maximize utility at state $s$. A coordinate set $I$ is sufficient when any two states that agree on $I$ also agree on $\Opt$. The optimizer quotient object then collapses exactly those states with the same optimal-action set, yielding the canonical abstraction through which every decision-preserving abstraction factors.

Section~\ref{sec:foundations} makes that collapse boundary explicit: any surjective abstraction either factors through the optimizer quotient object or erases a decision-relevant distinction, and any attempt to treat such extra erasure as physically feasible runs into the same no-collapse obstruction used later in the hardness chain. In \textbf{Set}, this quotient object is the familiar coimage of $\Opt$, canonically equivalent to $\operatorname{range}(\Opt)$.\LH{QT5}\LH{AB2}\LH{AB3}\LH{AB4}

The results in this section are organized as one derivation chain rather than a theorem ledger. The opening results establish the structural core from counting, set theory, arithmetic, and logic; later results then place the same core in complexity-theoretic and physical terms. Theorems 1--14 are derived from pure mathematics, while Theorems 15--17 add explicit empirical lower-bound premises. One irreducible empirical claim (EC1) supplies the thermodynamic scale: irreversible bit operations carry a positive lower-bound cost, with Landauer furnishing a universal floor. Counting and logic provide the structural part of the chain. Two former claims (EC2: finite resources, EC3: finite speed) are now derived from first principles.

\subsubsection{1. Counting Gap Theorem}

We begin with the discrete counting fact that every later information and complexity bound depends on: bounded total capacity and positive integer cost imply a finite observation budget.

\begin{theorem}[Counting Gap]\label{thm:counting-gap-intro}
Let $\varepsilon, C \in \mathbb{N}$ with $\varepsilon > 0$ and $C > 0$. If each information-acquisition event consumes $\varepsilon$ cost units and the total available capacity is $C$, then for every $N \in \mathbb{N}$:
\[
  \varepsilon \cdot N \leq C \;\Rightarrow\; N \leq C.
\]
Equivalently, $c_{\max}=C$ is a finite upper bound on the number of possible events. No bounded system with positive integer event cost can perform infinitely many events.
\LH{BA10}
\end{theorem}

\begin{proof}
This is a discrete theorem: $\varepsilon > 0$ in $\mathbb{N}$ implies $1 \leq \varepsilon$, so each event costs at least one unit. Therefore $N = 1\cdot N \leq \varepsilon \cdot N \leq C$.
\end{proof}

\noindent\textbf{Consequence.} Continuous abstractions require arbitrarily fine state discrimination. Bounded computational systems cannot do this. Discrete models are forced by the mathematics.

\subsubsection{2. Bayesian Optimality}

Once observations are finite, the next question is how to update beliefs optimally from those observations.

\begin{theorem}[Bayes Minimizes Expected Log Loss]\label{thm:bayes-optimal}
For distributions $p, q$ over hypothesis space $H$, cross-entropy satisfies:
\[
  \mathrm{CE}(p, q) = H(p) + D_{\mathrm{KL}}(p \| q)
\]
where $D_{\mathrm{KL}}(p \| q) \geq 0$ (Gibbs' inequality, from $\log x \leq x - 1$).
Therefore $q = p$ uniquely minimizes $\mathrm{CE}(p, q)$.
\LH{FN7}\LH{FN12}\LH{FN14}
\end{theorem}

\begin{proof}
Gibbs' inequality: $\log x \leq x - 1$ implies $\sum p \log(p/q) \geq 0$. Equality iff $p = q$.
\end{proof}

\noindent\textbf{Consequence.} Bayesian updating is optimal. The proof uses only $\log x \leq x - 1$ and Lean's type theory.

\subsubsection{3. Measure Theory Typing}

To state those updates correctly, we next distinguish the type of a quantitative claim from the stronger normalization required for a probabilistic claim.

\begin{theorem}[Measure-Theoretic Typing]\label{thm:measure-prerequisite}
\begin{enumerate}
\item Every quantitative claim carries an explicit measure: $(\mu, f)$ where $f$ is measurable.
\item Every stochastic claim additionally requires $\mu$ to be a probability measure ($\mu(\Omega) = 1$).
\item Counting measure on $\{0,1\}$ has mass $2$, not $1$: not a probability measure.
\end{enumerate}
\LH{MN1}\LH{MN2}\LH{MN10}\LH{MN11}
\end{theorem}

\begin{proof}
(1) ``$\mathbb{E}[f]$'' is undefined without specifying which measure integrates $f$. The expression $\int f\,d\mu$ requires $\mu$.
(2) $P(A) = \mu(A)/\mu(\Omega)$ requires $\mu(\Omega) < \infty$; normalization $P(\Omega) = 1$ requires $\mu(\Omega) = 1$.
(3) Counting measure: $\mu(\{0\}) = 1$, $\mu(\{1\}) = 1$, so $\mu(\{0,1\}) = 2 \neq 1$. To normalize, divide by 2.
\end{proof}

\noindent\textbf{Consequence.} Quantitative and stochastic claims are typed differently. Probability normalization is not automatic from measure structure.

\subsubsection{4. Universal Property of the Optimizer Quotient}

With the measurement layer typed, we can identify the canonical abstraction that preserves the decision boundary itself.

\begin{theorem}[Universal Property of the Optimizer Quotient]\label{thm:quotient-universal}
Let $Q = S/{\sim}$ be the optimizer quotient object where $s \sim s' \iff \Opt(s) = \Opt(s')$. For any surjective abstraction $\phi: S \to T$ that preserves the optimal action ($\phi(s) = \phi(s') \Rightarrow \Opt(s) = \Opt(s')$), there exists a unique factorization $\psi: T \to Q$ such that $\pi = \psi \circ \phi$ where $\pi: S \to Q$ is the quotient map.
\[
  \begin{array}{ccc}
    S & \xrightarrow{\pi} & Q \\
    \downarrow & & \uparrow \\
    T & \xrightarrow{\psi} & Q
  \end{array}
\]
The optimizer quotient object is the coarsest abstraction through which $\Opt$ factors. In \textbf{Set}, it is canonically equivalent to $\operatorname{range}(\Opt)$, i.e. the familiar coimage/image factorization of the optimizer map.
\LH{QT1}\LH{QT2}\LH{QT3}\LH{QT5}\LH{QT7}
\end{theorem}

\begin{proof}
Since $\phi$ preserves $\Opt$, we have $\phi(s) = \phi(s') \Rightarrow \Opt(s) = \Opt(s')$. By surjectivity, for each $t \in T$ pick some $s$ with $\phi(s) = t$ and define $\psi(t) = \pi(s)$. This is well-defined: if $\phi(s) = \phi(s') = t$, then $\Opt(s) = \Opt(s')$, so $\pi(s) = \pi(s')$. Uniqueness follows because $\pi = \psi \circ \phi$ forces $\psi(t) = \pi(s)$ for any $s$ with $\phi(s) = t$.
\end{proof}

\noindent\textbf{Consequence.} The optimizer quotient object is canonical. In \textbf{Set}, the familiar construction is not exotic: it is the coimage of $\Opt$, canonically identified with the image/range of $\Opt$. Any other state abstraction that preserves optimal actions must refine it.

\subsubsection{5. Witness-Checking Duality}

Once the canonical abstraction is fixed, we can ask how many state comparisons any sound verifier must inspect to certify that no counterexample exists.

\begin{theorem}[Checking-Witnessing Duality]\label{thm:checking-duality}
For boolean decision problems with $n$ coordinates, let $W(n) = 2^{n-1}$ be the witness budget and $T(n)$ be the checking budget. Any sound checker satisfies:
\[
  T(n) \geq W(n) = 2^{n-1}
\]
No finite checker with fewer than $2^{n-1}$ witness pairs can be sound on all problem instances.
\LH{WD1}\LH{WD2}\LH{WD3}
\end{theorem}

\begin{proof}
$\emptyset$-sufficiency fails iff $\exists s, s'$ with $\Opt(s) \neq \Opt(s')$. There are $2^n$ states, so $\binom{2^n}{2} = 2^{n-1}(2^n - 1)$ pairs. A sound refutation must cover all pairs that could witness failure. For each pair family $P$ with $|P| < 2^{n-1}$, there exists an $\Opt$ function whose unique failure pair is not in $P$. Therefore $|P| \geq 2^{n-1}$ is necessary for soundness.
\end{proof}

\noindent\textbf{Consequence.} Verification has a lower bound. A checker must examine exponentially many witness pairs to be sound. This is a counting lower bound on the information required for verification.

\subsubsection{6. Bayes from Counting}

That verification bound rests on a simpler structural fact: probability, Bayes, and the normalized quotient score already emerge from counting on finite sets.

\begin{theorem}[Probability Axioms from Counting]\label{thm:bayes-from-counting}
The chain from counting to Bayesian updating and the normalized quotient score:
\begin{enumerate}
\item $P(A) = |A|/|\Omega|$ satisfies Kolmogorov axioms: nonnegativity, normalization, and additivity for disjoint sets.
\item $P(H|E) = P(E|H) \cdot P(H) / P(E)$ follows from the definition of conditional probability.
\item Conditioning reduces entropy: $H(H|E) \leq H(H)$, from nonnegativity of KL divergence.
\item $\DQ = I(H;E)/H(H) = 1 - H(H|E)/H(H)$ lies in $[0,1]$.
\end{enumerate}
\LH{BC1}\LH{BC2}\LH{BC3}\LH{BC4}\LH{BC5}
\end{theorem}

\begin{proof}
(1) $|A| \geq 0$ (counting), $|\Omega|/|\Omega| = 1$ (normalization), $|A \cup B| = |A| + |B|$ for $A \cap B = \emptyset$ (additivity).
(2) $P(H|E) := P(H \cap E)/P(E)$. Rearranging: $P(H|E) \cdot P(E) = P(H \cap E) = P(E|H) \cdot P(H)$.
(3) $H(H|E) = H(H) - I(H;E)$ and $I(H;E) = D_{\mathrm{KL}}(P_{H,E} \| P_H \otimes P_E) \geq 0$ by Gibbs.
(4) $I(H;E) \leq H(H)$ by (3), so $\DQ = I(H;E)/H(H) \in [0,1]$.
\end{proof}

\noindent\textbf{Consequence.} Bayesian updating is derived from counting measure. Probability theory, Bayes' theorem, information theory, and the normalized quotient score form a single chain from counting elements of a finite set.

\subsubsection{7. Fisher Rank = Structural Rank}

With that probabilistic chain in place, the next results compare independent mathematical formalisms and ask whether they recover the same structural invariant. In the finite-coordinate model, $\mathrm{srank}$ is just the cardinality of the relevant-coordinate support, so the question is whether those formalisms recover that same support-size quantity.\LH{SK1}\LH{SK2}\LH{SK3}

\begin{theorem}[Fisher Information = Structural Rank]\label{thm:fisher-rank-srank}
Define the Fisher information score of coordinate $i$ as:
\[
  \mathrm{score}(i) = \begin{cases} 1 & \text{if } i \text{ is relevant} \\ 0 & \text{otherwise} \end{cases}
\]
Then:
\[
  \sum_{i=1}^{n} \mathrm{score}(i) = \mathrm{srank}(\mathcal{D})
\]
The diagonal Fisher information matrix $I(\theta)_{ii} = \mathrm{score}(i)$ has rank $\mathrm{srank}(\mathcal{D})$.
\LH{FS1}\LH{FS2}
\end{theorem}

\begin{proof}
$\sum_i \mathrm{score}(i) = \sum_i \mathbf{1}[\text{isRelevant}(i)] = |\{i \mid \text{isRelevant}(i)\}| = \mathrm{srank}$.
For the diagonal matrix: $\mathrm{rank} = |\{i \mid \mathrm{score}(i) \neq 0\}| = |\{i \mid \text{isRelevant}(i)\}| = \mathrm{srank}$.
\end{proof}

\noindent\textbf{Consequence.} The decision manifold has dimension exactly the size of the relevant-coordinate support. The Fisher information sees the same coordinates already counted by $\mathrm{srank}$. The Cramer-Rao bound follows: any estimator has difficulty $\geq 1/\mathrm{srank}$.

\subsubsection{8. Entropy-Rank Inequality}

The first bridge is information-theoretic: quotient entropy cannot exceed the number of coordinates that actually matter.

\begin{theorem}[Entropy Bounded by Structural Rank]\label{thm:entropy-rank}
For a decision problem $\mathcal{D}$ with boolean coordinate spaces ($X_i = \{0,1\}$):
\[
  H(\mathcal{D}) \leq \mathrm{srank}(\mathcal{D})
\]
where $H(\mathcal{D}) = \log_2(\mathrm{numOptClasses})$ is the Shannon entropy of the optimizer quotient object.
\LH{IT3}\LH{IT4}
\end{theorem}

\begin{proof}
The relevant coordinate set $R$ is sufficient (Theorem~\ref{thm:resolution-sufficient}): $\mathrm{Opt}$ factors through $\pi_R : S \to \{0,1\}^{|R|}$. Therefore $\mathrm{numOptClasses} \leq |\{0,1\}^{|R|}| = 2^{\mathrm{srank}}$. Therefore $H(\mathcal{D}) = \log_2(\mathrm{numOptClasses}) \leq \log_2(2^{\mathrm{srank}}) = \mathrm{srank}$.
\end{proof}

\noindent\textbf{Consequence.} The optimizer quotient object carries at most $\mathrm{srank}$ bits of decision-relevant information. Coordinates outside the relevant set are information-theoretically inert.

\subsubsection{9--14. Later Bridge Theorems}

The next six introductory results are bridge theorems. They are proved in their dedicated sections, but we state them here in compressed form because they all serve the same purpose: independent frameworks recover the same structural core once the optimizer quotient object is fixed.

\begin{theorem}[Thermodynamic Cost from Counting]\label{thm:fi-coincide}
For functional set $F \subseteq \Omega$, the counting expression $-\log_2(|F|/|\Omega|)$ gives the bit cost, and a Landauer floor converts that bit count into a universal energy lower bound. In the sharpened physical interface, explicit nonnegative mismatch and residual terms are added above that floor to represent constrained-process dissipation. The mismatch branch is now justified as far as the existing mathematics allows: distinct strictly positive finite input distributions yield strictly positive KL mismatch, which is then rounded upward into the discrete lower-bound units used by the thermodynamic accounting and injected into the Wolpert decomposition. A second derived branch now captures the strongest residual theorem supported by the current finite machinery: for finite discrete computational-state processes, positive forward edge flow together with decision-relevant edge asymmetry yields strict residual positivity by an exhaustive local split on the reverse edge. If the reverse flow is positive, the pairwise KL branch applies. If the reverse flow vanishes, the process performs an irreversible one-way state transition, and the existing Landauer-scaled transition-cost theorem supplies a positive lower-bound witness. That local argument is now also packaged as a process-level theorem: any finite computational-state process with at least one such asymmetric forward edge already carries a theorem-level residual witness. The broader stopping-time / absolute-irreversibility residual regime remains a separate imported premise. The thermodynamic statement is therefore the counting statement plus a floor-plus-overhead empirical conversion law; a derived mismatch branch, a derived finite residual branch from this exhaustive edge split, or the cited broader stopping-time branch can each force strict separation above the Landauer floor, and any declared structural lower bound absorbed by mismatch only strengthens the resulting energy estimate.
\LH{FI3}\LH{FI6}\LH{FI7}\LH{WC1}\LH{WC2}\LH{WC3}\LH{WM1}\LH{WM3}\LH{WM4}\LH{WM5}\LH{WM6}\LH{WR1}\LH{WR2}\LH{WR3}\LH{WR4}\LH{WR5}\LH{WR6}\LH{WR7}\LH{WR8}\LH{WR9}\LH{WR10}\LH{WP1}\LH{WP5}\LH{WP6}\LH{WP7}\LH{WP8}
\end{theorem}

\begin{proof}
The mathematical part is only the counting step: finite measure gives the bit count $-\log_2(|F|/|\Omega|)$. Landauer then supplies the universal floor for the energy-per-bit conversion. The bridge is therefore structural counting plus one declared empirical lower-bound premise.
\end{proof}

\noindent\textbf{Consequence.} The thermodynamic lower bound is not a separate combinatorial structure; it is the counting bound expressed in physical units.

\begin{theorem}[Wasserstein Distance = Structural Rank]\label{thm:wasserstein-bridge}
Optimal-transport cost grows with the number of distinguishable futures and therefore tracks the same relevant-coordinate complexity measured by $\mathrm{srank}$.
\LH{W1}\LH{W2}\LH{W3}\LH{W4}
\end{theorem}

\begin{proof}
When all states lead to one future, the diagonal coupling gives zero transport cost. Once multiple futures must be distinguished, off-diagonal mass incurs positive cost. The same branching structure counted by $\mathrm{srank}$ therefore reappears as transport cost.
\end{proof}

\noindent\textbf{Consequence.} Optimal transport recovers the same complexity signal as the coordinate-relevance analysis.

\begin{theorem}[Rate-Distortion = Structural Rank]\label{thm:rate-distortion-bridge}
At zero distortion, the minimum lossless rate needed to preserve optimal actions is exactly $\mathrm{srank}$, i.e. exactly the cardinality of the relevant-coordinate support, and compression below that support size necessarily introduces decision errors.
\LH{RD1}\LH{RD2}\LH{RD3}\LH{RS1}\LH{RS2}\LH{RS3}\LH{RS4}\LH{RS5}
\end{theorem}

\begin{proof}
Zero distortion means that decision-equivalent states may be merged, but decision-distinguishing states may not. The quotient therefore partitions the source into exactly the classes that must remain distinguishable. Their count determines the minimum lossless rate, which is the same structural quantity measured by $\mathrm{srank}$.
\end{proof}

\noindent\textbf{Consequence.} Description length and relevant-coordinate support size coincide at zero distortion.

\begin{theorem}[Information Requires Nontriviality]\label{thm:nontriviality-counting}
If decision-relevant information exists, then the state space cannot be trivial: the existence of information forces at least two distinguishable states.
\LH{FP1}\LH{FP2}\LH{FP3}\LH{FP4}\LH{FP5}\LH{FP6}\LH{FP7}
\end{theorem}

\begin{proof}
If the state space had at most one element, then every map out of it would be constant, including $\Opt$. A constant decision boundary has zero entropy and carries no distinguishable information. Taking the contrapositive gives the stated nontriviality requirement.
\end{proof}

\noindent\textbf{Consequence.} Nontriviality is derived from the existence of information rather than postulated as an extra axiom.

\begin{theorem}[Second Law Structure from Counting]\label{thm:second-law-counting}
Maintaining ordered behavior requires continuous distinction between correct and incorrect transitions; in this sense, the structural content of the Second Law follows from counting and verification, while the numerical thermodynamic scale is added later.
\LH{FP8}\LH{FP9}\LH{FP10}\LH{FP11}\LH{FP12}\LH{FP13}\LH{FP14}
\end{theorem}

\begin{proof}
If there are more non-target than target continuations, unguided evolution overwhelmingly drifts away from the target set. Staying ordered therefore requires repeated discrimination between correct and incorrect continuations. That discrimination is an information-acquisition burden, and the thermodynamic lower-bound premise enters only after that structural burden is fixed.
\end{proof}

\noindent\textbf{Consequence.} The order-versus-verification tradeoff is established before any temperature-dependent conversion is introduced.

\begin{theorem}[Erasure Increases Environment Entropy]\label{thm:landauer-structure}
Whenever a state reduction deletes distinctions, those distinctions must reappear in environmental degrees of freedom, so the structural content of Landauer's principle is a transfer statement and only the numerical thermodynamic scale is empirical.
\LH{FP15}
\end{theorem}

\begin{proof}
Reducing $n$ distinguishable states to fewer than $n$ outputs removes distinctions from the local description. Those distinctions cannot vanish from the bookkeeping; they must be displaced into the environment. The structural theorem is therefore a conservation-of-distinguishability statement, with temperature only setting the energetic scale of that transfer.
\end{proof}

\noindent\textbf{Consequence.} Landauer's principle separates cleanly into a mathematical transfer structure and an empirical lower-bound energy scale.

\subsection{Irreducible Empirical Claims}\label{sec:irreducible-empirical}

One claim is irreducibly empirical. Two others (finite resources, finite speed) are derived from first principles.

\begin{definition}[Irreducible Empirical Claim EC1]\label{def:empirical-claims}
\textbf{EC1 (Temperature):} Entropy has energy cost. There exists $\varepsilon > 0$ such that each irreversible bit operation costs at least $\varepsilon$ joules, with $\varepsilon \geq k_B T \ln 2$. Landauer gives the universal floor. A separate floor-plus-overhead interface treats any strictly larger per-bit dissipation as an explicit additional premise rather than an informal interpretation.
\end{definition}

\begin{theorem}[EC2 Derived from Uncertainty]\label{thm:ec2-derived}
Finite resources (EC2) follows from the existence of uncertainty:
\begin{enumerate}
\item Uncertainty exists: $\exists X, 0 < P(X) < 1$ (observation).
\item No bounds $\Rightarrow$ infinite information acquisition $\Rightarrow$ know everything $\Rightarrow$ $P(X) \in \{0,1\}$ for all $X$.
\item Contrapositive: uncertainty exists $\Rightarrow$ bounds exist.
\end{enumerate}
\LH{IA7}\LH{IA9}\LH{IA11}\LH{IA12}\LH{IA13}
\end{theorem}

\begin{proof}
Let observer have capacity $C$ and access $A \leq C$. System has entropy $H$. If $C \geq H$, observer can access all information (IA8). Full access eliminates uncertainty (IA9). Contrapositive: uncertainty implies $A < H$. If observer uses full capacity ($A = C$), then $C < H$ (IA11). Therefore capacity is bounded relative to system entropy. This is EC2.
\end{proof}

\begin{theorem}[EC3 Derived from Non-Contradiction]\label{thm:ec3-derived}
Finite signal speed (EC3) follows from the law of non-contradiction:
\begin{enumerate}
\item A computational step = state change: before $\neq$ after (definition).
\item At any instant, state is unique (function property).
\item If step takes $t = 0$: before-state and after-state at same instant.
\item Same instant $\Rightarrow$ state = before AND state = after.
\item But before $\neq$ after $\Rightarrow$ contradiction (violates non-contradiction).
\item Therefore $t_{\text{before}} \neq t_{\text{after}}$, so $\Delta t > 0$.
\item Information propagation = sequence of steps, each taking $\Delta t > 0$.
\item Speed = distance / time. Time $> 0$ $\Rightarrow$ speed $< \infty$.
\end{enumerate}
\LH{FPT4}\LH{FPT5}\LH{FPT6}\LH{FPT8}\LH{FPT10}
\end{theorem}

\begin{proof}
A step changes state from $A$ to $B$ where $A \neq B$. The timeline is a function $T \to S$. If $t_{\text{before}} = t_{\text{after}}$, then $\text{timeline}(t) = A$ and $\text{timeline}(t) = B$ for the same $t$. By function uniqueness, $A = B$. Contradiction. Therefore $t_{\text{before}} \neq t_{\text{after}}$. With ordering, $\Delta t > 0$ (FPT5). Information propagation requires steps; each takes positive time; total time is positive; speed is finite. QED from pure logic.
\end{proof}

\noindent\textbf{Consequence.} EC2 and EC3 are not axioms. EC2 follows from uncertainty existing. EC3 follows from non-contradiction. \textbf{One irreducible empirical claim remains: EC1 (temperature).} The empirical content is that irreversible bit operations have a positive lower-bound cost, with Landauer giving a universal floor and stronger constrained-process bounds allowed above it; the existence of a cost is structural.

\subsubsection{15. Assumption Necessity}

\begin{theorem}[Physical Claims Need Empirically Justified Assumptions]\label{thm:assumption-necessity}
If a physical claim $\phi$ is:
\begin{enumerate}
\item Provable from assumption set $\Gamma$ (sound formal system)
\item Falsifiable (some model $m$ with $\neg \mathrm{Holds}(m, \phi)$)
\item Derived from empirically disciplined $\Gamma$ (all axioms empirically justified)
\end{enumerate}
Then $\Gamma$ contains at least one physical axiom that is empirically justified.
\LH{AN3}\LH{AN4}
\end{theorem}

\begin{proof}
Suppose $\Gamma$ contains only mathematical axioms (no physical axioms). Mathematical axioms hold in all models (by soundness). Therefore $\phi$ holds in all models. But (2) says $\phi$ fails in some model $m$. Contradiction. Therefore $\Gamma$ contains at least one non-mathematical axiom. By (3), this axiom must be empirically justified.
\end{proof}

\noindent\textbf{Consequence.} Physical axioms cannot be hidden. Physical consequences require declared and empirically justified physical assumptions.

\subsubsection{16. Energy Bound}

\begin{theorem}[Energy Lower Bound]\label{thm:energy-information}
For any physical system resolving decision problem $\mathcal{D}$ at temperature $T > 0$:
\[
  E_{\mathrm{decision}} \geq k_B T \cdot H_{\mathrm{nats}}(\mathcal{D})
\]
where $H_{\mathrm{nats}}(\mathcal{D}) = \ln(\mathrm{numOptClasses})$ is the natural-log Shannon entropy of the optimizer quotient object.
\LH{EI1}
\end{theorem}

\begin{proof}
Three steps:
\begin{enumerate}
\item $\mathrm{srank}$ bits are required to resolve $\mathcal{D}$. Each bit requires one state transition.
\item Each transition costs at least the Landauer floor $k_B T \ln 2$. In the stronger constrained-process model, explicit nonnegative mismatch and residual dissipation terms may raise that per-bit lower bound further. The mismatch branch is no longer merely postulated: if the actual and designed finite distributions differ, the existing KL machinery yields a strictly positive mismatch witness, which is then converted into the discrete lower-bound units of the thermodynamic model. A second finite residual branch is also now derived by an exhaustive edge split on finite computational-state processes: if the reverse edge flow is positive, asymmetric forward/reverse edge flows yield strictly positive two-point KL divergence; if the reverse edge flow is zero, the process performs an irreversible one-way state transition and the existing Landauer-scaled transition-cost theorem yields a positive discrete lower-bound term. The broader stopping-time branch remains an imported premise. Therefore $E \geq \mathrm{srank} \cdot k_B T \ln 2$ remains the universal floor, and any declared constrained-process contribution tightens it monotonically.
\item $H_{\mathrm{nats}} = \ln(\mathrm{numOptClasses}) \leq \mathrm{srank} \cdot \ln 2$ (Theorem~\ref{thm:entropy-rank}).
\end{enumerate}
Combining: $E \geq \mathrm{srank} \cdot k_B T \ln 2 \geq k_B T \cdot H_{\mathrm{nats}}$.
\end{proof}

\noindent\textbf{Consequence.} Thermodynamic cost tracks information content: energy per decision $\geq k_B T$ per nat of decision entropy. This is derived from the Landauer floor (Theorem 9) composed with the entropy-rank inequality (Theorem 8); in the explicit decomposed constrained-process interface, a theorem-level KL mismatch witness, a theorem-level finite residual witness produced by the exhaustive reverse-edge split, its process-level witness packaging, or the cited broader stopping-time residual branch only strengthens the estimate.\LH{WC1}\LH{WC4}\LH{WM3}\LH{WM4}\LH{WM6}\LH{WR6}\LH{WR7}\LH{WR8}\LH{WR9}\LH{WR10}\LH{WP1}\LH{WP5}\LH{WP6}

\subsubsection{17. Thermodynamic Uncertainty Relations}

\begin{theorem}[TUR Bridge (Conditional)]\label{thm:tur-bridge}
Conditional on stochastic thermodynamics (Barato--Seifert 2015):
\begin{enumerate}
\item Transition probabilities are normalized: $p_{ij} \geq 0$, $\sum_j p_{ij} = 1$.
\item The Second Law enforces non-negative entropy production: $\sigma \geq 0$.
\item The TUR bound constrains precision: $\mathrm{Var}(J)/\langle J \rangle^2 \geq 2/\sigma$.
\item Multiple futures require positive entropy production.
\end{enumerate}
Precision requires entropy dissipation, scaling with structural rank.
\LH{TUR1}\LH{TUR2}\LH{TUR5}\LH{TUR6}
\end{theorem}

\begin{proof}
(1) Standard probability. (2) The Second Law (empirical assumption). (3) TUR is a theorem of stochastic thermodynamics: for any current $J$ in a Markov process with entropy production $\sigma$, precision is bounded by $\mathrm{Var}(J)/\langle J \rangle^2 \geq 2/\sigma$.
(4) If $\mathrm{srank} > 1$, distinguishing $k > 1$ futures requires information acquisition. By the TUR bound, finite precision requires $\sigma > 0$. The minimal $\sigma$ scales with $\log k$, hence with $\mathrm{srank}$.
\end{proof}

\noindent\textbf{Consequence.} Accepting fluctuation theorems yields an independent constraint: precision costs entropy production, with the bound scaling with srank. This is conditional on accepting the Second Law.

\subsection{Consistency with Existing Physical Theories}\label{sec:intro-justifies-physics}

The Counting Gap Theorem is consistent with and illuminates why discrete models work across physics:

\begin{itemize}
\item \textbf{Quantum Mechanics:} Discrete eigenstates are forced for bounded systems
\item \textbf{Statistical Mechanics:} Discrete microstates forced by finite accessible state count
\item \textbf{Thermodynamics:} Finite entropy forced by bounded information acquisition
\item \textbf{Lattice Gauge Theory:} Discretization respects the computational constraint
\end{itemize}

\noindent The paper rests on one irreducible empirical claim (Definition~\ref{def:empirical-claims}):

\begin{itemize}
\item \textbf{EC1 (Temperature):} $\varepsilon \geq k_B T \ln 2 > 0$ (entropy has a positive lower-bound energy cost)
\end{itemize}

\noindent EC2 (finite resources) is derived from the existence of uncertainty (Theorem~\ref{thm:ec2-derived}): if you don't know everything, your capacity is bounded. EC3 (finite speed) is derived from non-contradiction (Theorem~\ref{thm:ec3-derived}): a computational step must take positive time, or the state would be both $A$ and not-$A$ at the same instant. Everything else---nontriviality, second law structure, Landauer structure, EC2, EC3---follows from counting, logic, or observation, with the numerical thermodynamic scale entering only through the declared lower-bound premise in EC1.

\subsection{Our Formal Tools: Coordinate Sufficiency}\label{sec:intro-formal-tools}

We study decision problems with coordinate structure. The finite coordinate spaces $X_i$ are forced by the Counting Gap. Boolean coordinates $X_i = \{0,1\}$ are elementary: one physical transition per bit.

A coordinate set $I$ is sufficient if knowing $s_I$ determines the optimal action:
\[
s_I = s'_I \implies \Opt(s) = \Opt(s')
\]

We prove complexity classifications for three meta-problems:

\begin{itemize}
\item \textbf{SUFFICIENCY-CHECK:} Is $I$ sufficient? coNP-complete
\item \textbf{MINIMUM-SUFFICIENT-SET:} Find minimal $I$. coNP-complete
\item \textbf{ANCHOR-SUFFICIENCY:} Exists sufficient $I$ with $|I| \leq k$? $\Sigma_2^P$-complete
\end{itemize}

Stochastic variants are PP-complete. Sequential variants are PSPACE-complete. Strict hierarchy: static $\subset$ stochastic $\subset$ sequential.

\subsection{Machine Verification}\label{sec:intro-verification}

The Lean 4 formalization consists of:
\begin{itemize}
\item 24,347 lines of Lean code
\item 1,082 theorems
\item 0 \texttt{sorry} placeholders
\item 0 hidden axioms
\end{itemize}

\paragraph{What is machine-verified.}
Bayesian optimality: proved from $\log x \leq x - 1$ alone. Thermodynamic consequences: conditional on a Landauer floor or any stricter declared per-bit lower bound. \textbf{Reduction correctness}: the membership-preservation property of each reduction (e.g., $\varphi$ is a tautology $\Leftrightarrow$ $I = \emptyset$ is sufficient) is fully machine-verified.

\paragraph{What is a standard claim.}
\textbf{Polynomial-time computability} of the reduction functions is a standard claim stated as an explicit hypothesis in the Lean formalization. Full Turing-machine verification of polynomial-time bounds remains an open problem in formal verification---no formalization project has achieved end-to-end polynomial-time TM proofs for concrete reductions. The gap is made visible by representing polytime as a hypothesis rather than hiding it behind definitions that conflate size bounds with time bounds. Machine-generated assumption ledger records all dependencies.

\subsection{Paper Structure}\label{sec:intro-structure}

Section~\ref{sec:foundations}: formal definitions (decision problems, sufficiency, DQ).
Section~\ref{sec:physical-grounding}: physical grounding (Counting Gap, three laws).
Section~\ref{sec:hardness}: hardness proofs (coNP, $\Sigma_2^P$, PP, PSPACE).
Section~\ref{sec:dichotomy}: regime separation.
Section~\ref{sec:tractable}: six tractable subcases.
Sections~\ref{sec:engineering-justification}--\ref{sec:simplicity-tax}: engineering corollaries.
Appendix: Lean listings.

\paragraph{Claim-stamp notation.}
Claims are stamped with theorem references $[T:n; P:m]$ and Lean handles (\LH{H1}).

\begingroup
\footnotesize
\noindent
\begin{tabularx}{\linewidth}{@{}>{\raggedright\arraybackslash}p{0.12\linewidth}>{\raggedright\arraybackslash}p{0.40\linewidth}>{\raggedright\arraybackslash}p{0.43\linewidth}@{}}
\toprule
\textbf{Regime} & \textbf{Question} & \textbf{Placement} \\
\midrule
Static & SUFFICIENCY / MINIMUM / ANCHOR & coNP / coNP / $\Sigma_2^P$ \\
Stochastic & STOCHASTIC-SUFFICIENCY & PP-complete \\
Sequential & SEQUENTIAL-SUFFICIENCY & PSPACE-complete \\
\bottomrule
\end{tabularx}
\endgroup

\leanmeta{\LHrng{QT}{1}{3}, \LHrng{WD}{1}{3}, \LHrng{BC}{1}{5}, \LH{CCC1}, \LHrng{CCC}{5}{6}, \LHrng{DC}{1}{2}, \LHrng{DC}{9}{10}, \LH{FN14}, \LH{MN5}, \LH{FI7}, \LHrng{W}{1}{4}, \LHrng{RD}{1}{3}, \LHrng{RS}{1}{5}, \LHrng{TUR}{1}{2}, \LHrng{TUR}{5}{6}.}
